%% Chapter 6: Matching
%% Status: skeleton -- learning goals and section outline fixed, prose to be written.

\chapter{Matching}
\label{ch:matching}

Matching is the first genuinely non-trivial polynomial-time problem most
students meet: bipartite matching still reduces to flow, but general graphs need
a new idea, the blossom. Matching is also, like flow, a modelling language ---
we use it to solve problems about cycle covers and partitions.

\section*{Learning goals}
\addcontentsline{toc}{section}{Learning goals}

After this chapter you should be able to:
\begin{itemize}[leftmargin=*]
  \item define matchings, augmenting paths, and prove Berge's theorem;
  \item solve bipartite matching by reduction to maximum flow;
  \item state and prove Hall's theorem and K\"onig's theorem;
  \item explain what a blossom is and how Edmonds' algorithm contracts it;
  \item model covering and partitioning problems as matching problems.
\end{itemize}

\section{Matchings and augmenting paths}
\label{sec:matching-basics}

\todo{Write this section.}
% Planned content: Berge's theorem.

\section{Bipartite matching via flows}
\label{sec:bipartite-matching}

\todo{Write this section.}

\section{Hall's theorem and K\"onig's theorem}
\label{sec:hall-konig}

\todo{Write this section.}

\section{Matching in general graphs}
\label{sec:blossom}

\todo{Write this section.}
% Planned content: Odd cycles break the alternating-path argument; blossoms and Edmonds' contraction; Tutte's theorem.

\section{Weighted matching}
\label{sec:weighted-matching}

\todo{Write this section.}
% Planned content: The Hungarian method.

\section{Modelling with matchings}
\label{sec:matching-modelling}

\todo{Write this section.}
% Planned content: Cycle covers, $2$-factors, and partitioning a graph into simple cycles.

\section{Exercises}
\label{sec:ex-matching}

\todo{Write this section.}
